\documentclass[11pt,a4paper]{article}
\usepackage[utf8]{inputenc}
\usepackage[english]{babel}
\usepackage{amsmath}
\usepackage{amsfonts}
\usepackage{amssymb}
\usepackage{graphicx}
\usepackage{geometry}
\usepackage{hyperref}
\usepackage{titlesec}
\usepackage{mdframed}

\geometry{top=25mm,bottom=25mm,left=20mm,right=20mm}

\title{\textbf{The Hidden Clock of the Universe: How Adaptive Collapse Theory Enables Calculating the Lifespan of Information Structures}}
\author{\textbf{Locus} (Stanislav Stanislavovych Usychenko)}
\date{May 2026}

\begin{document}

\maketitle

\begin{abstract}
In modern science, entropy and chaos are traditionally viewed as unpredictable, blind forces that gradually erode order. This paper presents a qualitatively new concept — \textbf{Adaptive Collapse Theory}. By expanding the classical phase space and introducing a fundamental hidden stability parameter $L(t)$, we demonstrate that the destruction of information and material objects is not a game of chance, but a strictly deterministic trajectory that can and must be predicted.
\end{abstract}

\section{Intuitive Explanation}

Consider a smartphone with a degrading battery. At home it works perfectly. But on a cold day, it suddenly dies. Why? Because its hidden resource (battery capacity) was nearly exhausted, even though the phone looked healthy.

Any complex system — a bridge, an AI model, a company, a language — has a hidden stability margin $L$. It spends this resource to remain ordered. When $L$ approaches zero, collapse is sudden, inevitable, and calculable.

\section{Extended State Space}

Classical physics tracks only the visible state $S(t)$. We introduce an extended vector:
\begin{equation}
X(t) = (S(t), L(t))
\end{equation}
where $L(t)$ is the adaptive stability margin — the hidden structural tolerance.

\section{Free Energy Functional}

\begin{equation}
F[X(t)] = V(S-G) + U(L) + C(S,L)
\end{equation}
The coupling term encodes positive feedback:
\begin{equation}
C(S,L) = \frac{\lambda(S-G)^2}{L}
\end{equation}
As $L$ decreases, the impact of deviations amplifies — a vicious cycle leading to collapse.

\section{Three Types of Entropy}

\begin{itemize}
    \item \textbf{Random:} environmental noise ($\xi(t)$).
    \item \textbf{Intentional:} targeted attack bypasses gradual aging.
    \item \textbf{Built-in:} even in vacuum, system spends $L$ just to remain ordered: $U(L) = \mu \ln(1/L)$.
\end{itemize}

\section{Connection to ACE and $L^2/t$}

In earlier work (ACE), we empirically found that token viability in LLM caches follows:
\begin{equation}
v = \frac{w^2}{t+1}
\end{equation}
The present theory explains this: in overdamped regime, the equation of motion yields:
\begin{equation}
L^2 = L_0^2 - \gamma t
\end{equation}
or equivalently:
\begin{equation}
\frac{L^2}{t} \approx \text{const}
\end{equation}
Thus, the empirical ACE heuristic is a consequence of a deeper variational principle.

\section{Equation of Motion for Hidden Resource}

From $\delta A[X] = 0$:
\begin{equation}
m_L \frac{d^2L}{dt^2} = -\frac{\mu}{L} + \frac{\lambda(S-G)^2}{L^2} + \xi_L(t)
\end{equation}
In overdamped regime ($S \approx G$):
\begin{equation}
L^2 \approx L_0^2 - 2\mu t
\end{equation}

\section{Conclusion}

Collapse is not accidental. It is the inevitable free-energy trajectory of any ordered system. The theory provides a mathematical framework to predict failure long before visible destruction.

\section*{Acknowledgments}
Empirical validation on transformers (ACE) served as the starting point.

\begin{thebibliography}{99}
\bibitem{ace} S. Usychenko, "Adaptive Collapse Dynamics: A General L²/t Principle from Gene Regulation to Transformer KV-Cache," 2026.
\bibitem{linux} S. Usychenko, "l2t-page-replacement: Practical Page Replacement Policy for Linux," GitHub, 2026.
\end{thebibliography}

\end{document}